\section{Related Work}

\subsection{Palmprint Recognition}

Palmprint recognition has historically relied on local texture structure, including line, orientation, and filter-bank representations. Recent survey evidence describes the field's transition toward deep feature learning, with learned representations increasingly used for ROI processing, feature extraction, matching, and robustness analysis \cite{gao2025deeplearning}. This paper follows the embedding-learning direction rather than presenting a new ROI detector or a handcrafted texture descriptor.

Gabor and other texture filters remain important context for palmprint recognition because they encode local ridge and crease patterns. However, the current project evidence does not support a Gabor-centered paper claim: the fixed-Gabor B2 variant is a useful baseline, but B6 provides the stronger final result on the Tongji multi-seed cross-session protocol. Accordingly, this work does not claim that Gabor filtering is the source of the final improvement.

\subsection{Metric Learning for Biometric Embeddings}

Biometric recognition systems depend on embedding geometry: genuine pairs should remain close while impostor pairs should be separated, especially at strict false-acceptance operating points. Supervised contrastive learning explicitly pulls same-class examples together and pushes different-class examples apart \cite{khosla2020supervised}. The B1 baseline in this project uses this idea together with cross entropy, forming a strong baseline for the Tongji cross-session protocol.

Margin-based normalized classifiers provide another route to discriminative embeddings. ArcFace adds an angular margin to normalized features and class weights, encouraging separation in angular space \cite{deng2019arcface}. BNNeck was introduced in person re-identification as a normalization neck that separates feature spaces used for metric comparison and classification supervision \cite{luo2019strong}. B6 adapts these two ideas to the palmprint setting: ArcFace shapes the angular classifier geometry, while BNNeck provides the post-BN embedding used for cosine verification.

\subsection{Cross-Session and Cross-Domain Robustness}

Cross-session and cross-domain robustness are central issues for deployable biometric systems. The report reviewed for this draft mentions several palmprint domain-adaptation, style-transfer, and synthetic-data works, but many of those entries still need full author-title-venue-year verification before they are added to the bibliography. For this reason, the present draft uses them only as motivation for future literature expansion, not as cited evidence.

This paper focuses on a narrower and directly measured setting: Tongji cross-session recognition with S1->S2 and S2->S1 protocols. The results should therefore be read as evidence for cross-session Tongji verification, not as a broad cross-domain benchmark claim. External datasets and verified cross-domain comparisons remain future work.
